% ============================================================
% 第2章 问题分析
% ============================================================

\section{问题分析}

\subsection{问题一分析}
针对问题一，板凳龙沿螺距为$55~\text{cm}$的等距螺线盘入，各把手中心均位于螺线上，需在龙头前把手恒速$1~\text{m/s}$条件下确定$0\sim300~\text{s}$内所有把手的位置和速度。由于所有把手均在同一螺线上，不宜将各把手坐标作为独立变量，而应以\textbf{螺线极角$\theta$}为位置参量进行\textbf{参数化}建模。首先，将龙头运动转化为弧长参数化问题，利用\textbf{阿基米德螺线弧长微分公式}建立\textbf{弧长}函数，由龙头速度恒为$1~\text{m/s}$得$\boxed{S(\theta_1(t))=t}$，通过求解该一元方程得到\textbf{龙头极角}$\theta_1(t)$，从而确定任意时刻龙头前把手对应的螺线参数及其位置。随后，利用相邻把手距离固定的\textbf{几何约束}$\boxed{\left| P(\theta_{i+1})-P(\theta_i) \right| = L_i}$，构造关于$\theta_{i+1}$的非线性方程，并利用\textbf{余弦定理$+$二分检验搜索求解}，之后逐节递推求得所有把手的极角。最后，对龙头恒速条件直接求导得到\textbf{龙头角速度}，其余把手速度则通过先对$x$,$y$方向上分量分别求解，再合成为总速度的方式计算。我们利用“弧长参数化+逐节递推+二分法+隐式求导”的策略，使多体运动问题分解为一系列一维问题，有效降低了计算复杂度。



\subsection{问题二分析}
针对问题二，我们需要在已知板凳龙所有把手的运动状态的条件下，去判断舞龙队在盘入过程中何时首次发生板凳碰撞。已知每节板凳具有一定的长度和宽度，因此，本文首先将每节板凳抽象为以相邻把手连线为中心线、长度为$341cm$或$220cm$，宽度为$30cm$的矩形。其次，对于任意时刻$t$，利用问题一中的递推模型计算全部把手位置；随后由相邻把手确定各节板凳的位置和姿态，并利用$SAT$矩形碰撞判定方法判断任意两节非相邻板凳是否发生重叠。最后，利用和问题一中类似的基于二分法的时间搜索，求解首次检测碰撞时刻，将该时刻定义为舞龙队的临界盘入时刻$t^*$。

\subsection{问题三分析}
对第三个问题的扩展性和复杂点进行分析。